\documentclass[11pt,a4paper]{article}

% Lingua e caratteri (compilare con XeLaTeX)
\usepackage{fontspec}
\usepackage[italian]{babel}

% Matematica, figure e collegamenti
\usepackage{amsmath, amssymb}
\usepackage{graphicx}
\usepackage[hidelinks]{hyperref}

\title{Titolo dell’articolo}
\author{Nome Cognome}
\date{\today}

\begin{document}

\maketitle

\begin{abstract}
  Breve riassunto del lavoro: il problema, l’approccio e i risultati principali.
\end{abstract}

\section{Introduzione}
\label{sec:introduzione}

Questo è un articolo di esempio. Le lettere accentate (à, è, é, ì, ò, ù) funzionano
direttamente grazie a XeLaTeX. Per maggiori dettagli, si veda la
sezione~\ref{sec:sviluppo} e il riferimento~\cite{knuth1984}.

\section{Sviluppo}
\label{sec:sviluppo}

Un’equazione numerata:
\begin{equation}
  \label{eq:euler}
  e^{i\pi} + 1 = 0
\end{equation}

e un sistema allineato:
\begin{align}
  \nabla \cdot \vec{E} &= \frac{\rho}{\varepsilon_0}, \\
  \nabla \cdot \vec{B} &= 0.
\end{align}

L’equazione~\eqref{eq:euler} lega cinque costanti fondamentali.

\section{Conclusioni}

Riassumi qui le conclusioni principali.

\bibliographystyle{plain}
\bibliography{bibliografia}

\end{document}
